\chapter{Conclusion}
\label{ch:conclusion}

The premise of this project was that using a large language model to generate fuzzing inputs directly, rather than to synthesise a generator that emits them, would yield a more diverse seed corpus and better fuzzing performance. The idea was to remove a layer of indirection, and utilise the LLM's learned knowledge of the input format to write those inputs directly, in a syntactic and semnatically aware manner. We implemented \toolname{}, which embodies this direct generation approach, and compared it against the original generator-synthesising ELFuzz on two systems under test: jsoncpp and cpython3 (Chapter~\ref{ch:evaluation}). Prior work supports the broader idea: LLMs have been used to generate test inputs directly under a simple crash oracle with effective results \cite{xia_2023_universal,lin_2025_lmfuzz}, and ELFuzz has shown that coupling an LLM with coverage feedback to evolve input \emph{generators} is highly effective \cite{chen_2025_elfuzz}. We combined these threads in \toolname{}, which uses an LLM to generate seed inputs directly and then evolves them under coverage guidance, and we hypothesised that this more direct route would yield a more diverse corpus and, in turn, better fuzzing performance.

On the evidence gathered, that hypothesis is not supported. Across the two systems under test we evaluated, the original generator-synthesising ELFuzz matched or beat \toolname{} on every end-to-end measure (Chapter~\ref{ch:evaluation}). Its seed corpora reached more coverage on their own ($568$ versus $163$ edges on jsoncpp, $19{,}634$ versus $9{,}790$ on cpython3), and the edge sets they covered almost entirely subsumed those of \toolname{}, rather than a complementary part of the target. Downstream, the gap was erased on the small jsoncpp - where AFL++ saturated the reachable frontier within minutes regardless of the seed - but it persisted on the larger cpython3, where ELFuzz stayed ahead for the whole campaign. We did not observe the diversity advantage that motivated the approach.

Several factors plausibly contribute, and they bound how strongly this result should be read. The most immediate is the generator itself. Hardware constraints limited us to a small local model, Qwen2.5-Coder-1.5B \cite{hui_2024_qwen25coder}, where the original ELFuzz experiments used the far larger CodeLlama-13B \cite{rozire_2024_code,chen_2025_elfuzz}. The cost of this was visible directly in the inputs: not one of the jsoncpp pool's inputs was well-formed JSON, and only about a third of the cpython3 pool compiled as valid Python. A model that cannot reliably emit a complete, valid instance of a simple format will struggle to produce the structurally varied inputs that the hypothesis claims will be more effective. The indirection of synthesising a generator may in fact be what shields ELFuzz from this weakness - a generator can encode the format's structure once and emit many valid instances, whereas direct generation must get each input right on its own.

A second factor is the evolution rule. \toolname{} curates its pool purely on coverage, retaining an input only when it contributes a new edge. This is a strict criterion, and it interacts poorly with the way ELFuzz appears to win: its \texttt{produce} step samples ~1 million candidate inputs from the evolved generators, and that sheer quantity - not just per-input quality - is a large part of where its coverage comes from. A evolution strategy that entirely optimises for coverage growth (\toolname{}'s jsoncpp corpus minimised to just three inputs) will discard the redundant-but-varied volume that benefits the downstream mutation fuzzer. Optimising for coverage alone may therefore be the wrong objective for a corpus that is about to be handed to AFL++.

\section{Future Work}

These are early results, based on a single run, a small model, and only two systems under test (Chapter~\ref{ch:evaluation}), so there is plenty of room to build on them. The most pressing limitation is the model itself, and it showed in the outputs: very few of the generated inputs were even valid, so most of them only exercised the target's parser and never reached the logic behind it. The natural next step is to repeat the experiment with a larger, more capable model on hardware that can host it, and to run the campaigns for longer, such as the 24 hours used in the original ELFuzz evaluation \cite{chen_2025_elfuzz} rather than the six we used. That would show whether the validity problem is fundamental to direct generation or simply an artefact of model size, and give any real difference in seed quality more time to show. Input validity is worth tackling head-on as well, since getting the model to produce valid inputs is the key to exercising more than just the parser. Constraining the model so that it can only emit tokens that are valid in the target format, or using a model trained on that format, are both promising ways to do this.

It would also help to test on more targets, and on richer ones. jsoncpp turned out to be too simple to be informative, as AFL++ exhausted its reachable code almost immediately for both tools, so the quality of the seeds made no difference downstream. Larger, more complex targets such as programming-language interpreters and compilers are a much better test of the idea, because their input space is nowhere near exhausted within a realistic time budget. The way \toolname{} curates its pool is also worth revisiting. It keeps an input only when it adds new coverage, but part of ELFuzz's strength seems to come from sheer quantity, producing around a million inputs before minimising. A more permissive rule that keeps a larger, more varied pool, or that rewards structural and semantic diversity directly rather than coverage alone, may suit a downstream mutation fuzzer better than the strict consolidation we used here.
